\section{引言}
\label{Sec:Introduction}

区块链是一种去中心化的数字账本技术，能够在多台计算机之间安全地记录和验证交易，从而保障数据完整性和透明性~\cite{bitcoin}。作为区块链生态系统中最具代表性的代币，比特币（Bitcoin，BTC）在 2024 年 11 月达到了 1.8 万亿美元的显著市值，体现了该技术的影响力与成功~\cite{coinmarketcap}。


在区块链技术的基础上，去中心化应用（Decentralized Applications，DApps）已成为一项关键创新~\cite{buterin2014next}。这些应用运行在部署于以太坊（Ethereum）、Polygon、Binance Smart Chain（BSC）等不同区块链网络上的智能合约之上。DApp 的扩展反映了应用开发向去中心化模式转变的趋势，并提供了去中心化金融（Decentralized Finance，DeFi）和 GameFi 等多样化服务。加密货币钱包、区块链浏览器等广泛基础设施支撑着这一持续演进的生态系统，并共同推动了区块链领域的发展活力~\cite{tang2022blockchain}。

智能合约事件和交易日志对于追踪区块链应用中的操作至关重要，它们支撑着跨链桥、钱包以及 NFT 市场等场景中的多种功能。伪造这些日志会削弱系统安全性，导致资产盗窃、欺诈以及互操作性问题，进而威胁用户信任。尽管已有研究考察了特定领域中的漏洞，例如跨链桥~\cite{Xscope,liao2024smartaxe,wang2024xguard}、NFT 所有权~\cite{sleepminting}、代码不一致性~\cite{doccon}以及代币行为~\cite{TokenScope}，但目前仍缺乏面向多类应用的事件伪造综合研究。本文通过引入\emph{幻影事件（Phantom Events）}这一新型漏洞类别来弥补上述空白，该类别揭示了区块链中\emph{日志伪造}的广泛风险。不同于以往针对孤立场景的研究，本文系统总结了去中心化应用中的事件误用现象，并探索了用于识别幻影事件的不同检测方法。在已有研究中，XScope~\cite{Xscope} 和 XGuard~\cite{wang2024xguard} 等检测工具在源代码层面运行，当仅有字节码可用时，其检测能力会受到限制。相反，SmartAxe~\cite{liao2024smartaxe} 和 TokenScope~\cite{TokenScope} 虽然分析字节码，但仅依赖字节码的方法会遗漏事件参数细节。Guidi 等人~\cite{sleepminting}的方法仅关注交易日志，由于上下文分析有限，因而存在较高误报率。我们的工具结合了字节码、源代码（若可用）以及交易数据，能够跨多类漏洞进行综合检测，同时降低单层分析方法中常见的误报问题。

幻影事件是区块链系统中的一种隐蔽操纵，它会将正常的合约事件转化为执行未授权行为的隐藏通道。这类事件会破坏区块链交易中的信任基础，使攻击者能够绕过常规安全检查实施未授权活动。它们通过模仿或轻微篡改智能合约事件，误导事件监听器、用户界面以及区块链分析工具。此类漏洞利用会破坏数据完整性，并可能造成严重的经济损失和声誉损害。例如，2024 年 2 月 9 日，攻击者通过伪造转账事件造成了 104 万美元的损失~\cite{zero_transfer}。

本文研究重点在于理解并归类不同类型的幻影事件。我们对近期区块链安全报告和文章进行了全面调研，以收集利用幻影事件相关漏洞的真实攻击案例。此外，我们审计了多个活跃智能合约，并构建了一个用于处理区块链事件的安全模型（\cref{sec:threat}），该模型是分析和归类所收集攻击场景的基础（\cref{sec:vector}）。进一步地，我们开发了一个名为 \tool 的工具，用于检测与幻影事件相关的漏洞（\cref{sec:detection}）。\tool 的设计面临多项挑战：（1）许多合约仅以字节码形式存在，缺乏事件语义，从而增加了分析难度；（2）幻影事件会模仿合法事件的格式，因此难以区分；（3）判断是否存在异常还需要理解合约的业务逻辑。

利用所提出的安全模型和检测工具，我们在多种区块链应用中识别出了大量漏洞和潜在风险（详见\cref{sec:evaluation}）。结果表明，幻影事件漏洞具有广泛分布性，并影响区块链生态系统中的多类应用。通过审计工作，我们发现了此前未被报告的不一致日志记录漏洞实例，在这些实例中，链上问题会导致链下记录出现差异。此外，我们的工具还识别出了此前未被检测到的链上历史事件问题，凸显了该工具在发现既往事件方面的有效性。在合约级分析中，我们的工具优于现有解决方案，在检测幻影事件漏洞方面表现出更高的准确性。进一步地，我们在加密货币钱包、区块链浏览器和跨链桥等多个真实应用中识别出了安全问题，并将这些发现报告给项目团队，通过负责任披露申请漏洞赏金。最后，我们在\cref{sec:mitigation}中提出了应对这些漏洞的缓解策略。

简而言之，本文的主要贡献如下：
\begin{itemize}
\item \textbf{新颖的攻击分类。}
我们较早地对区块链交易日志伪造相关攻击进行了系统分析和分类。本文提出了一个能够刻画五类攻击的安全模型。我们还识别出了此前未被报告的不一致日志记录漏洞，其中链上问题会导致链下记录出现差异。
\item \textbf{高效的检测方法。}
我们开发了一种检测机制，在识别幻影事件方面超过了现有工具，并成功发现了此前未被报告的攻击交易。
\item \textbf{实践意义。}
我们在区块链浏览器、加密货币钱包、GameFi、DeFi 和 NFT 市场等多种区块链应用中识别出了真实问题。值得注意的是，我们在一个跨链中继器中发现了漏洞，而该中继器也被其他项目使用，其中包括一个市值超过 2.5 亿美元的项目。截至目前，我们已向项目团队和漏洞赏金平台报告了六个加密货币钱包问题、一个区块链浏览器问题以及一个跨链桥问题，其中五个已被确认、一个已被修复，并通过负责任披露获得了总计 600 美元的赏金。
\end{itemize}
